%-----------------------------------------------------------------
% Rodrigo Chang, 2019
% Banco de Guatemala
% Programa de Estudios Superiores 2019-2020
% Diapositivas del curso "Programación I"
%
% Descargo de responsabilidad.
%   El contenido de este documento es responsabilidad exclusiva de
% su(s) autores. Cualquier opinión expresada no necesariamente
% representa la postura institucional del Banco de Guatemala. 
% 
% Contenido.
%   El contenido de estas diapositivas está construido a partir de 
% modificaciones al material provisto por Ana Bell, Eric Grimson y 
% John Guttag para el curso 6.0001 Introduction to Computer Science 
% and Programming in Python, otoño 2016, Massachusetts Institute of 
% Technology: MIT OpenCourseWare, https://ocw.mit.edu
%
% Licencia: 
% Creative Commons BY-NC-SA
% https://creativecommons.org/licenses/by-nc-sa/4.0/.
%
% Formato. 
%   El formato utilizado para la elaboración de este documento 
% ha sido modificado a partir del repositorio: 
%       https://github.com/rinnert/statlearning-python-pwf
%
% Kurt Rinnert, 2019
%   Copying and distribution of this file, with or without 
% modification, are permitted in any medium without royalty provided 
% the copyright notice and this notice are preserved. This file is 
% offered as-is, without any warranty.
%
%-----------------------------------------------------------------
\documentclass[mathserif]{beamer}
% Para configurar las diapositivas en relación 16:9, 
% agregar la opción aspectratio=169

% Título de la presentación y tópicos a tratar
\usepackage{arev, arevmath}
\usepackage[scaled]{cabin}
%\usepackage[T1]{fontenc}
\usepackage[super]{nth}
\usepackage{pifont}
\usepackage{wasysym}
\usepackage{tabularx}
\usepackage{array}
\usepackage{booktabs}
\usepackage{boldline}
\usepackage{colortbl}
\usepackage{amsmath}
\usepackage{bm}
\usepackage{tcolorbox}
\usepackage{adjustbox}
\usepackage{minibox}
\usepackage{makecell}
\usepackage{adjustbox}
\usepackage{textcomp}
\usepackage[absolute,overlay]{textpos}
\setlength{\TPHorizModule}{1mm}%
\setlength{\TPVertModule}{1mm}%
\tcbuselibrary{skins}
\usepackage{graphicx}
\usepackage[spanish, es-tabla]{babel}
\usepackage[utf8]{inputenc}
\usepackage{csquotes}

\usepackage{comment}
\usepackage{svg}

\usepackage{verbatim}
\usepackage{pdfpages}

% Columnas múltiples
\usepackage{multicol}

% -------------------------------------------
%	LISTINGS
% -------------------------------------------
% Paquetes necesarios para los códigos
\usepackage{listings}
\usepackage{color}
% Configuración de listings (pendiente)
%\input{conf/listingStyles}
% Redefinición de titulos de caption
\renewcommand\lstlistingname{\textbf{Código}}
\renewcommand\lstlistlistingname{Códigos}

% Definición de colores para Python (Jupyter notebooks)
\definecolor{python_keyword_color}{RGB}{0, 128, 0}
\definecolor{python_string_color}{RGB}{186, 34, 34}
\definecolor{python_comment_color}{RGB}{1, 121, 121}
\definecolor{frame_numbers_color}{rgb}{0.5,0.5,0.5}
\definecolor{python_identifier_color}{RGB}{0,0,0}
\definecolor{jupyter_background}{RGB}{247, 247, 247}
\definecolor{jupyter_frame}{RGB}{207, 207, 207}

% Definicion de estilo para Python
\lstdefinestyle{custom_python}{
	language=python,			% the language of the code
	backgroundcolor=\color{jupyter_background}, 
	% the size of the fonts that are used for the code
	basicstyle=\linespread{1}\small\ttfamily,	
	% Configuración de colores
	commentstyle=\color{python_comment_color},	% comment style
	identifierstyle=\color{python_identifier_color},
	keywordstyle=\mdseries\color{python_keyword_color},	% keyword style
	numberstyle=\tiny\color{frame_numbers_color},	
	stringstyle=\color{python_string_color}, % string literal style 
	captionpos=t,				% sets the caption-position to bottom
	breakatwhitespace=false,	% sets if automatic breaks should only happen at whitespace
	breaklines=true,			% sets automatic line breaking
	morekeywords={},			% if you want to add more keywords to the set
	deletekeywords={eps},		% delete keywords from the given language
	escapeinside={\%*}{*)},		% if you want to add LaTeX within your code
	framesep=3pt,
	keepspaces=true,			% keeps spaces in text,       
	numbers=none, 	% (none, left, right)
	numbersep=5pt,	% how far the line-numbers are from the code
	% if not set, the frame-color may be changed on line-breaks within not-black text (e.g. comments (green here))
	rulecolor=\color{jupyter_frame}, 
	% show spaces everywhere adding particular underscores; it overrides 'showstringspaces'
	showspaces=false,                
	% underline spaces within strings only
	showstringspaces=false,          
	% show tabs within strings adding particular underscores
	showtabs=false,                  
	stepnumber=0,	% the step between two line-numbers      
	% sets default tabsize to 2 spaces
	tabsize=2,                       
	% lets you use non-ASCII characters; for 8-bits encodings only, does not work with UTF-8
	extendedchars=true,      
	% Support additional characters            
	literate=
  {á}{{\'a}}1 {é}{{\'e}}1 {í}{{\'i}}1 {ó}{{\'o}}1 {ú}{{\'u}}1
  {Á}{{\'A}}1 {É}{{\'E}}1 {Í}{{\'I}}1 {Ó}{{\'O}}1 {Ú}{{\'U}}1
  {à}{{\`a}}1 {è}{{\`e}}1 {ì}{{\`i}}1 {ò}{{\`o}}1 {ù}{{\`u}}1
  {À}{{\`A}}1 {È}{{\'E}}1 {Ì}{{\`I}}1 {Ò}{{\`O}}1 {Ù}{{\`U}}1
  {ä}{{\"a}}1 {ë}{{\"e}}1 {ï}{{\"i}}1 {ö}{{\"o}}1 {ü}{{\"u}}1
  {Ä}{{\"A}}1 {Ë}{{\"E}}1 {Ï}{{\"I}}1 {Ö}{{\"O}}1 {Ü}{{\"U}}1
  {â}{{\^a}}1 {ê}{{\^e}}1 {î}{{\^i}}1 {ô}{{\^o}}1 {û}{{\^u}}1
  {Â}{{\^A}}1 {Ê}{{\^E}}1 {Î}{{\^I}}1 {Ô}{{\^O}}1 {Û}{{\^U}}1
  {œ}{{\oe}}1 {Œ}{{\OE}}1 {æ}{{\ae}}1 {Æ}{{\AE}}1 {ß}{{\ss}}1
  {ű}{{\H{u}}}1 {Ű}{{\H{U}}}1 {ő}{{\H{o}}}1 {Ő}{{\H{O}}}1
  {ç}{{\c c}}1 {Ç}{{\c C}}1 {ø}{{\o}}1 {å}{{\r a}}1 {Å}{{\r A}}1
  {€}{{\EUR}}1 {£}{{\pounds}}1,   
}

% Configurar el listing para el documento actual
\lstset{
	style = custom_python, 
	frame = single
}

% -------------------------------------------
%	OTRAS CONFIGURACIONES
% -------------------------------------------
\makeatletter
\newcommand{\antsize}{\@setfontsize{\antsize}{4pt}{4pt}}
\makeatother
\newcommand{\at}{\makeatletter @\makeatother}

\newcommand{\cmark}{\ding{51}}%
\newcommand{\bottomline}[1]{\vskip0pt plus 1fill{\alert{#1}\phantom{g}\vskip 1.0mm}}

% Rutas para las imágenes
%\graphicspath{{figs/}{figs/common/}}

% Separación de filas en las tablas más amplia
\renewcommand{\arraystretch}{1.2}


% -------------------------------------------
%	ESTILO DE LA PRESENTACIÓN
% -------------------------------------------

\mode<presentation>%
{
	\usetheme{default}
	%\usetheme[width=2.5cm]{PaloAlto}
	\usecolortheme{dove}
	\useoutertheme{infolines}
	% oder auch nicht

	% BANGUAT Colors
	\definecolor{BANGUATblue}{RGB}{37,86,162} % 0x255682
	\definecolor{BANGUATorange}{RGB}{255,130,0} % 0xff8200
	\definecolor{BANGUATgreen}{RGB}{0,100,0}
	\definecolor{BANGUATdark}{RGB}{80,80,80} % 0x505050
	\definecolor{BANGUATlight}{RGB}{120,120,120}
	\definecolor{BANGUATbrown}{RGB}{178,91,0}

	\definecolor{codebg}{rgb}{0.95,0.95,0.95}

	% Color theme
	\setbeamercolor{alerted text}{fg=BANGUATorange}
	\setbeamercolor{frametitle}{fg=BANGUATblue}
	\setbeamercolor{title}{fg=BANGUATblue}
	\setbeamercolor{subtitle}{fg=BANGUATorange}
	\setbeamercolor{normal text}{fg=BANGUATdark}
	\setbeamercolor{author in foot}{fg=BANGUATblue}
	\setbeamercolor{item}{fg=BANGUATblue}
	\setbeamercolor{footline}{fg=BANGUATblue}
	%\setbeamercolor{item projected}{bg=BANGUATorange}
	%\setbeamercolor{item projected}{fg=white}

	\setbeamertemplate{headline}
	{}
	\setbeamertemplate{frametitle}
	{
		%\textbf{{\insertframetitle\phantom{g}}}\\
		%\textbf{\insertframetitle\phantom{g}}\\
		\textbf{\underline{\insertframetitle\phantom{g}}}\\
		%\textbf{\underline{\insertframetitle}}\\
		\vskip 1.0mm
		%{\color{UOLgold}\hrule height 2pt}
		%\par
	}
	\addtobeamertemplate{frametitle}{}{\vspace{-1em}}
	\setbeamertemplate{footline}{
		{%
			\textbf{ \hskip 3.0mm\insertshorttitle\phantom{.}---\phantom{.}\insertshortinstitute\hfill\insertframenumber\,/\,\inserttotalframenumber\hskip 3.0mm} 
		}
	}

	\setbeamertemplate{navigation symbols}{}%remove navigation symbols
	\setbeamertemplate{itemize items}[circle]
	\setbeamertemplate{enumerate items}[fg=BANGUATblue]
	\setbeamercolor{itemize items}{fg=BANGUATblue}
	\setbeamercolor{sidebar}{bg=BANGUATblue}
	\setbeamercolor{title in sidebar}{fg=BANGUATorange}
	\setbeamercolor{author in sidebar}{fg=BANGUATorange}
	\setbeamercolor{section in sidebar}{fg=BANGUATorange}
}



% -------------------------------------------
%	CONFIGURACIÓN DE DIAPOSITIVA DE TÍTULO
% -------------------------------------------

% Título del curso
\title[Programación I]{\textbf{%
Introducción a las ciencias de la computación\\%
{\it y programación en Python}}\\%
%\href{https://ocw.mit.edu/courses/electrical-engineering-and-computer-science/6-0001-introduction-to-computer-science-and-programming-in-python-fall-2016/}{\tiny\it Basado en el curso ``Introduction to Computer Science and Programming in Python'' por Ana Bell, Eric Grimson, y John Guttag.}%
\vspace{3em}%
}
%\vspace{-2.5cm}{}

% Autor y correo electrónico
\author{\href{mailto:?to=Rodrigo Chang <rrcp@banguat.gob.gt>&subject=PES Programación I}{Rodrigo Chang}}

% Información del instituto
\institute[\href{https://www.banguat.gob.gt/}{Banco de Guatemala}]%
{\color{BANGUATblue} \bfseries \href{https://www.banguat.gob.gt/}{Banco de Guatemala}\\%
\vspace{3mm}%
\includegraphics[width=0.13\textwidth]{figs/common/logo-banguat.png}}

% Sin fecha en particular
\date{}

% Información de licencias y copyright
\titlegraphic{
	 %\texorpdfstring{\vspace{-3cm}}{}
	 \vspace{-1.5cm}
	 \begin{minipage}[b][1cm][b]{0.25\textwidth}%
	    \color{BANGUATlight}\antsize
		%Copyright \textcopyright~2019\\
		\href{mailto:?to=Rodrigo Chang <rrcp@banguat.gob.gt>&subject=PES Programación I}{Rodrigo Chang <rrcp@banguat.gob.gt>}\\
		Este material está construido a partir de modificaciones al material provisto por Ana Bell, Eric Grimson y John Guttag para el curso \textit{6.0001 Introduction to Computer Science and Programming in Python}, otoño 2016, Massachusetts Institute of Technology: MIT OpenCourseWare, \href{https://ocw.mit.edu}{https://ocw.mit.edu}. Licencia: \href{https://creativecommons.org/licenses/by-nc-sa/4.0/}{Creative Commons BY-NC-SA}.
		
		%Copying and distribution of this file, with or without modification, are permitted in any medium without royalty provided the copyright notice and this notice are preserved.  This file is offered as-is, without any warranty.


		%Some of the figures in this presentation are taken from ``An Introduction to Statistical Learning, with applications in R''  (Springer, 2013) with permission from the authors: G. James, D. Witten,  T. Hastie and R. Tibshirani 
	 \end{minipage}%
	 \hspace{9cm}
}

% Restar uno al contador
\addtocounter{framenumber}{-1}

% Cajas de colores
\newcommand{\tabboxset}{\tcbset{enhanced, nobeforeafter, boxrule=0pt, boxsep=0pt, colback=codebg, colframe=codebg, coltext=BANGUATdark, rounded corners, arc=4pt, fonttitle={\bfseries\tiny}}}
\newcommand{\codeboxset}{\tcbset{enhanced, nobeforeafter, boxrule=0pt, boxsep=0pt, colback=codebg, colframe=codebg, coltext=BANGUATdark, rounded corners, arc=4pt, fonttitle={\bfseries\tiny}}}

% Comandos útiles
\newcommand{\orange}{\color{BANGUATorange}}
\newcommand{\blue}{\color{BANGUATblue}}
\newcommand{\dark}{\color{BANGUATdark}}
\newcommand{\R}{\mathbb{R}}
\newcommand{\dat}[1]{{\footnotesize\tt\orange #1}}
\newcommand{\mdat}[1]{\text{\dat{#1}}}
\newcommand{\val}[1]{{\footnotesize\tt\blue #1}}
\newcommand{\mval}[1]{\text{\val{#1}}}
\newcommand{\e}[1]{\emph{#1}}
\newcommand{\bh}{\hat{\beta}}
\newcommand{\h}{\hat}
\newcommand{\bl}{\bottomline}


\makeatletter
\newcommand{\includegraphicsdpi}[3]{%
	\pdfimageresolution=#1%
	\includegraphics[#2]{#3}%
	\pdfimageresolution=72%
}

\newcommand{\Quote}[2]{%
	\begin{center} 
		\begin{minipage}{0.7\textwidth} 
			\hrule
			\vskip 3mm
			\emph{{\color{BANGUATblue} ``#1''}\\
			{\color{BANGUATorange} --- #2}}
			\vskip 3mm
			\hrule
			\vskip 2mm
		\end{minipage}
	\end{center}
}

% Texto centrado en la diapositiva, quizás para el abstract o para indicar abajo de un quote
\newenvironment{blurb}%
	{\begin{center}\begin{minipage}{0.6\textwidth}\footnotesize}
	{\end{minipage}\end{center}}

% Environment para códigos o algoritmos. Parecido a un bottomline, se puede utilizar para dar definiciones o enfatizar en algún concepto
\newenvironment{cpage}%
	{\begin{center}\begin{minipage}{0.75\textwidth}}
	{\end{minipage}\end{center}}

% Bloque (similar a un pop-up) de alerta para indicar resultados en tablas, definiciones, fórmulas, etc.
% Primer argumento es el tamaño 
% Segundo argumento es el título del bloque
\newenvironment{popblock}[2]%
	{\begin{center}\begin{minipage}{#1}\footnotesize
		\begin{tcolorbox}[colframe=codebg, colback=white, colupper=BANGUATdark, title={\normalsize\bfseries\blue #2}]}
	{\end{tcolorbox}\end{minipage}\end{center}}
\makeatother

% Incluir diapositivas en PDF
% 1 - Título de la diapositiva
% 2 - páginas del PDF a incluir
% 3 - ruta hacia el archivo PDF
\newcommand{\includeslides}[3]{%
{%
\setbeamercolor{background canvas}{bg=}%
\addtocounter{framenumber}{1}%
\includepdf[pages={#1}, link=true, pagecommand={%
\Large\blue\bfseries\underline{%
#3\phantom{g}}%
\vspace{-1em}%
\vskip 1.0mm}]{#2}%
}}
\subtitle{\bfseries%
  {Control de flujo e iteraciones}\\%
  %{\tiny\it Control de flujo e iteraciones}\\%
}

% Inicio del documento
\begin{document}
\frame[plain]{
	%\vskip 1.0mm
	\titlepage
	%\vskip 1.0mm
}
%-----------------------------------------------------------------
% Empieza el contenido de la presentación
%-----------------------------------------------------------------

\begin{frame}{Abstract}
	\Quote{I taught myself how to program computers when I was a kid, bought my first computer when I was 10, and sold my first commercial program when I was 12.}{Elon Musk}
    
    % Everybody should learn to program a computer, because it teaches you how to think. - Steve Jobs
    
    % Whether you want to uncover the secrets of the universe, or you just want to pursue a career in the 21st century, basic computer programming is an essential skill to learn. - Stephen Hawking
    
    % I taught myself how to program computers when I was a kid, bought my first computer when I was 10, and sold my first commercial program when I was 12. - Elon Musk

	\begin{blurb}
		Veremos como obtener una entrada del usuario y generar una salida a través de la línea de comandos.
		
		Hablaremos de las condicionales y de las estructuras de iteración.
	\end{blurb}
\end{frame}

\section{Entradas y salidas}

\begin{frame}[fragile]{Cadenas de texto (introducción)}
    \begin{itemize}
        \item Son objetos de tipo {\blue no escalar} que permiten manejar {\orange números, letras, espacios y caracteres especiales}.
        \item Se encierran en {\blue comillas} dobles o simples.
\begin{lstlisting}
hello = "Hola mundo"
hi = 'Hola amigos del PES'        
\end{lstlisting}
        
        \item Podemos {\blue concatenar} cadenas de texto con el operador \lstinline|+|:
\begin{lstlisting}
nombre = "Antonio"
saludo = hello + " " + nombre
\end{lstlisting}
        
        \item Podemos efectuar otras operaciones: 
\begin{lstlisting}
molesto = hello + " " + nombre*3
\end{lstlisting}
    \end{itemize} 
\end{frame}


\begin{frame}[fragile]{Salida en consola}
    \begin{itemize}
        \item Utilizamos la función \lstinline|print|.
\begin{lstlisting}
print("Hola mundo")
\end{lstlisting}
        
        \item Podemos enviar {\blue varios argumentos} a la función para imprimirlos todos. Podemos separar los elementos con el caracter especificado en el argumento \lstinline|sep=' '|.
\begin{lstlisting}
pi=3.1415
print(pi, 2, "amigos")
print(1, 'Python Tricks', 'Dan Bader', sep=',')
\end{lstlisting}        

    \end{itemize}
\end{frame}

\begin{frame}[fragile]{Combinar variables en el texto}
    \begin{itemize}
        \item Será útil en muchas ocasiones, combinar en un mensaje de salida una o más variables.
        
        \item Para esto, utilizamos el operador \lstinline|%|. 
\begin{lstlisting}
pi=3.1415
print("Hola %s, Pi=%0.4f" % ('Rodrigo', pi))
\end{lstlisting}           
        
        \item También es posible utilizando diccionarios:
\begin{lstlisting}
print("Pi=%(pi)s" % {'pi' : pi})
\end{lstlisting}

        \item Otras especificaciones posibles: 
            \begin{itemize}
                \item \lstinline|%d| para números enteros.
                \item \lstinline|%s| para cadenas de texto.
                \item \lstinline|%f| para números de punto flotante.
            \end{itemize}
        
    \end{itemize}
    \bottomline{\url{https://realpython.com/python-string-formatting/}}
\end{frame}

\begin{frame}[fragile]{Entrada vía línea de comandos}
    \begin{itemize}
        \item Utilizamos el la función \lstinline|input("")|.
        
        \item Devuelve el valor escrito por el usuario a una variable:
\begin{lstlisting}
nombre = input("Ingresa tu nombre: ")
print("Hola " + nombre)
\end{lstlisting}

        \item \lstinline|input| devuelve objetos de tipo \lstinline|str|. Al trabajar con números se deben convertir los tipos:
\begin{lstlisting}
r = float(input("Ingresa el radio: "))
area = 3.14159*(r**2)
\end{lstlisting}

    \end{itemize}
\end{frame}

\section{Flujo de control}

\begin{frame}[fragile]{Condicionales}
    \begin{itemize}
        \item Podemos decidir si ejecutar un grupo de instrucciones con la siguiente estructura. \lstinline|<condition>| debe ser booleana.
\begin{lstlisting}
if <condition>:
    <expression>
    <expression>
    ...
\end{lstlisting}

        \item La cláusula \lstinline|else| indica las instrucciones a evaluar si la condicion es {\blue falsa}.
\begin{lstlisting}
if <condition>:
    <expression>
    ...
else:
    <expression>
    ...
\end{lstlisting}

    \end{itemize}
    
    \bottomline{La indentación es IMPORTANTE, esta es la forma en que Python estructura los bloques.}
\end{frame}

\begin{frame}[fragile]{Condicionales}
    \begin{itemize}
        \item La cláusula \lstinline|elif| permite añadir condiciones adicionales.
\begin{lstlisting}
if <condition>:
    <expression>
    ...
elif <condition>:
    <expression>
    ...
else:
    <expression>
    ...
\end{lstlisting}

    \end{itemize}
\end{frame}

\begin{frame}[fragile]{Bloques de código}
\begin{lstlisting}
x = float(input("Enter a number for x: "))
y = float(input("Enter a number for y: "))
if x == y:
   print("x and y are equal")
   if y != 0:
       print("therefore, x / y is", x/y)
elif x < y:
   print("x is smaller")
elif x > y:
   print("y is smaller")
print("thanks!")
\end{lstlisting}
\end{frame}


\section{Estructuras de iteración}

\begin{frame}[fragile]{Ciclo \lstinline|while|}
    \begin{itemize}
        \item Cuando queremos repetir una expresión {\blue hasta} que se cumpla una condición. 
\begin{lstlisting}
while <condition>:
    <expression>
    <expression>
    ...
\end{lstlisting}

        \item Nuevamente, \lstinline|<condition>| debe ser de tipo booleano.
        
        \item Si la condición es \lstinline|True|, el bloque se ejecuta. 
        
    \end{itemize}
\end{frame}

\begin{frame}[fragile]{Ciclo \lstinline|while|}
    \begin{itemize}
        \item Veamos un ejemplo con ciclo \lstinline|while|.
    \end{itemize}
    \lstinputlisting{codes/02-while-example.py}
\end{frame}

\begin{frame}[fragile]{Ciclos \lstinline|for|}
    \begin{itemize}
        \item Cuando queremos iterar sobre números de una secuencia.
        
\begin{lstlisting}
# more complicated with while loop
n = 0
while n < 5:
    print(n)
    n = n+1

# shortcut with for loop
for n in range(5):
    print(n)
\end{lstlisting}
    \end{itemize}
\end{frame}

\begin{frame}[fragile]{Ciclos \lstinline|for|}
    \begin{itemize}
        \item En cada iteración, \lstinline|<variable>| toma un valor.
\begin{lstlisting}
for <variable> in range(<some_num>):
    <expression>
    ...
\end{lstlisting}
        
        \item La primera vez, toma el valor más pequeño, \lstinline|0|. 
        \item La siguiente vez, toma el valor anterior + 1.

    \end{itemize}
\end{frame}

\begin{frame}[fragile]{Función \lstinline|range|}
    \begin{itemize}
        \item Los argumentos \lstinline|start| y \lstinline|step| son opcionales.
        \item Tienen valores por defecto: \lstinline|start = 0| y \lstinline|step = 1|.
        
        \item El ciclo llega hasta \lstinline|stop - 1|.
        %\item Devuelve un objeto de tipo \lstinline|'range'|.
    \end{itemize}
\begin{lstlisting}
# Sumamos números de 7 a 9
mysum = 0
for i in range(7, 10):
    mysum += i
    print(mysum)

# Sumamos números impares
mysum = 0
for i in range(5, 11, 2):
    mysum += i
    print(mysum)
\end{lstlisting}
    
    \bottomline{Note el uso del operador abreviado \lstinline|+=|.}
\end{frame}


\begin{frame}[fragile]{Sentencia \lstinline|break|}
    \begin{itemize}
        \item {\blue\bfseries Fuerza} la salida del ciclo actual.
        \item Instrucciones posteriores son ignoradas.
        \item Solamente actúa en el {\blue ciclo más interno}.
    \end{itemize}
\begin{lstlisting}
while <condition_1>:
    while <condition_2>:
        <expression_A>
        break
        <expression_B>
    <expression_C>
\end{lstlisting}
\end{frame}

\begin{frame}[fragile]{Sentencia \lstinline|break|}
\begin{lstlisting}
mysum= 0
for i in range(5, 11, 2):
    mysum += i
    if mysum == 5: 
        break
        mysum += 1 
print(mysum)
\end{lstlisting}

    \bottomline{¿Qué pasa en este programa?}
\end{frame}


\begin{frame}{Ciclos \lstinline|for| vs. \lstinline|while|}
    \begin{columns}
        \begin{column}{0.5\textwidth}
        \begin{cpage}
            \centering\lstinline|for|
        \end{cpage}
        \begin{itemize}
            \item Se {\blue conoce} el número de iteraciones.
            \item Admite \lstinline|break|.
            \item Utiliza un {\blue contador} o {\blue lista}.
            \item Se puede {\blue reescribir} un ciclo \lstinline|for| utilizando \lstinline|while|.
        \end{itemize}
        \end{column}
        \begin{column}{0.5\textwidth}
        \begin{cpage}
            \centering\lstinline|while|
        \end{cpage}
        \begin{itemize}
            \item Número de iteraciones no acotado. 
            \item Admite \lstinline|break|.
            \item Puede utilizar {\blue contador} pero {\orange debe inicializarlo y actualizarlo} en el ciclo.
            \item Puede {\orange no} ser posible {\blue reescribir} un ciclo \lstinline|while| utilizando \lstinline|for|.
        \end{itemize}
        \end{column}
    \end{columns}
\end{frame}
%-----------------------------------------------------------------
% Termina el contenido de la presentación
%-----------------------------------------------------------------
\end{document}